\section{Digital Twin and Financial Scenario Comparison}

The Financial Digital Twin is a deterministic what-if calculator that enables users to explore how changes in financial assumptions affect projected outcomes. It is an estimate, not a prediction of actual financial performance.

The Digital Twin uses the user's financial assessment as a baseline. The assessment inputs, including monthly income, expenses, accessible savings, EMI payments, the EMI inclusion flag, savings capacity, target amount, current goal savings and goal timeline, are reused directly without translation. The user may create a hypothetical scenario by modifying one or more of these inputs, for example increasing income, reducing expenses, increasing accessible savings, adding a one-time unexpected expense, or changing EMI payments.

The Twin recalculates the baseline and the scenario independently using the same deterministic formulas as the assessment engine. Monthly outflows add separately paid EMI to expenses; monthly surplus is income minus outflows; savings rate is surplus divided by income, null when income is zero; living-expense coverage is accessible savings divided by expenses, null when expenses are zero; mandatory-outflow coverage is accessible savings divided by outflows, null when outflows are zero; EMI burden is EMI divided by income, null when income is zero; post-shock savings is accessible savings minus a one-time unexpected expense, floored at zero; and unfunded shock is the shortfall when the expense exceeds savings.

Goal projection computes remaining amount, assumed monthly contribution, estimated months to goal, required contribution for a requested timeline, and whether the goal is on track. Estimated completion dates add projected whole calendar months to the as-of date and clamp to the last valid day of the target month, handling month-end boundaries such as 31 January plus one month yielding 28 or 29 February.

The comparison presents every numeric metric side by side with the baseline, showing the scenario value, the absolute difference, the percentage change where the baseline is nonzero and both values are available, and for rates an additional percentage-point difference. Goal comparison additionally shows remaining amount and assumed contribution differences, estimated-month difference, completion-date direction and whether goal achievement or on-track status changed. Recommendation comparison returns which recommendations are added, resolved, unchanged or severity-changed between baseline and scenario.

Each metric and recommendation carries an explanation with its label, formula, inputs, reason and assumptions. Heuristic recommendations include the configured threshold value with a note that the threshold is a heuristic. The interface displays an explicit estimate-not-prediction notice and the engine disclaimer, and never presents a scenario as an actual outcome.

The Digital Twin's educational value lies in helping users understand the sensitivity of their financial position to changes in income, expenses, savings and obligations, supporting more informed planning without guaranteeing future performance.